\section{Introduction}

An interactive system has to name the activity while the person is still
moving. A closed label such as ``sitting'' is too coarse for that response.
Sitting still, sitting and looking at a phone, and sitting and waving call
for different responses, yet a small label set collapses them into one
class. The system therefore describes the activity with a short English
phrase. A strong vision-language model can write that phrase, but running
it on every captured frame is too expensive for a phone, and uploading
every frame makes the user wait on the uplink.

The design splits the work. SmolVLM-256M-Instruct runs in the phone browser
and writes the phrase before any picture leaves the device
\cite{marafioti2025smolvlm}. The confidence of the phrase is the smallest
probability among the newly generated tokens. The JPEG stays on the phone
when this confidence is high and the phrase still agrees with the action
already on screen. Otherwise the phone uploads the last few frames. A
laptop on the same Wi-Fi network runs Qwen3.5-2B, writes its own phrase,
and the screen keeps or replaces the ongoing phrase according to the cosine
similarity of the two texts.

Two choices in the proposal have been changed because they fought that
goal. The proposal kept a fixed sentence bank and let CLIP retrieve one
bank sentence for the phone text and another for the picture. The screen
then showed the bank sentence. The running system keeps the model's own
phrase. The bank is used only when a free phrase must be scored against the
15 classes of the still-image evaluation. The proposal also planned a
compute-aware upload scheduler. That scheduler is the networking half of
the project. This midterm finishes the sentence-and-check path, including a
causal check over the last few frames, and states the scheduler work that
is still ahead.

The sentence path now has two extra gates in front of that upload. A live
frame is reduced to a \(32\times 32\) grayscale image. When an action is
already on screen, the leave count is zero, and the mean absolute
difference from the previous frame is below 10, the phone keeps the phrase
and does not run SmolVLM. A confident phrase that leaves the action is
held for one frame. The buffered frames go to Qwen only on a second
mismatch. A missing or low token probability does not wait. On a 60-image
sample the text-encoder matcher scores 38/60 for the phone model and 40/60
for Qwen. The same live rule, replayed on three frames from each test
clip, scores 1640/3783 on UCF101 and 660/1530 on HMDB51. A single Qwen
pass over those three frames scores 1785/3783 and 812/1530.

The rest of this report is organized as follows.
Section~\ref{sec:literature} places the prototype among activity
recognizers, vision-language models, and edge offloading systems.
Section~\ref{sec:objectives} states the objectives and which of them this
midterm meets.
Section~\ref{sec:design} gives the keep rule, the three-frame check, and
the upload control that is still to be built.
Section~\ref{sec:implementation} describes the browser client and the
laptop server.
Section~\ref{sec:evaluation} reports the still-image measurement and the
three-frame replay on UCF101 and HMDB51.
Section~\ref{sec:conclusion} states what the measurement supports and what
remains before the presentation on 20 November 2026.
